\begin{definition}[Names, interpretation, and extension]
Define the hierarchy of \(\mathbb P\)-names by ordinal recursion:
\[
  V^{\mathbb P}_\alpha
  \defeq
  \Pow\!\left(
    \left(\bigcup_{\beta<\alpha}V^{\mathbb P}_\beta\right)
    \times\mathbb P
  \right),
  \qquad
  V^{\mathbb P}\defeq
  \bigcup_{\alpha\in\operatorname{Ord}}V^{\mathbb P}_\alpha.
\]
A \(\mathbb P\)-name is a member of the class \(V^{\mathbb P}\).  Thus each
member of a name \(\nu\) is a pair \(\tuple{\mu,p}\), where \(\mu\) has lower
name rank and \(p\in\mathbb P\).  For a poset filter \(G\), rank recursion
defines
\[
  \nu^G
  \defeq
  \set{\mu^G:
        \exists p\in G\,(\tuple{\mu,p}\in\nu)}.
\]
For each \(x\in M\), the notation \(\check x\) is needed to exhibit a
\(\mathbb P\)-name in \(M\) whose interpretation is \(x\).  Define it by rank
recursion:
\[
  \check x\defeq
  \set{\tuple{\check y,1_{\mathbb P}}:y\in x}.
\]
Its standing is a set of name--condition pairs, hence a
\(\mathbb P\)-name, rather than the set \(x\) itself.  This distinction
is nonredundant because \(x\) need not have the pair structure required of a
name.  Because a nonempty poset filter contains \(1_{\mathbb P}\), induction on
the rank of \(x\) gives
\(\check x\in M\cap V^{\mathbb P}\) and \((\check x)^G=x\).  This calculation
is the exact link by which members of \(M\) enter \(M[G]\), and the
notation is used below in the name for the Cohen real.  It is not an appeal
to the cited forcing theorem below.
The associated extension is the set
\[
  M[G]
  \defeq
  \set{\nu^G:\nu\in M\cap V^{\mathbb P}}.
\]
Both recursions are well founded on ordinal rank; neither clause presupposes a
forcing relation.
\end{definition}

\begin{definition}[Semantic forcing under the fixed hypothesis]
\emph{Necessity.}
The qualifier \emph{semantic} is required to distinguish the exact external
relation defined here from the separate recursively defined forcing relation
whose construction is only cited below.

\emph{Definition.}
Let \(\varphi\) be a formula of the language of membership whose free variables
are among the pairwise distinct variables \(v_1,\ldots,v_n\), let
\(\nu_1,\ldots,\nu_n\in M\cap V^{\mathbb P}\), and let \(p\in\mathbb P\).
In this section the notation
\[
  p\forces^{\mathrm{sem}}_{M,\mathbb P}
    \varphi(\nu_1,\ldots,\nu_n)
\]
means exactly that, for every \(M\)-generic poset filter
\(G\subseteq\mathbb P\) with \(p\in G\) and every assignment \(s_G\) into
\(M[G]\) satisfying \(s_G(v_i)=\nu_i^G\) for \(1\leq i\leq n\),
\[
  (M[G],\in),s_G\entails\varphi.
\]
The preceding existence lemma prevents this universal condition from being
vacuous.

\emph{Standing.}
For the fixed \(M\) and \(\mathbb P\), semantic forcing is a metalinguistic
relation among a member \(p\) of \(\mathbb P\), a membership-language formula,
and a tuple of \(\mathbb P\)-names.

\emph{Nonredundancy.}
This relation is not ordinary satisfaction: it quantifies over every
\(M\)-generic poset filter containing \(p\).  It is also not the recursively
defined relation used inside \(M\) to prove definability and preservation;
no such recursive relation is defined in this manuscript.

\emph{Use.}
The notation is used immediately in the monotonicity lemma and later in the
source-pinned comparison with recursive forcing.
\end{definition}

